\section{Canonical Encodings}
\label{sec:encodings}

73-byte binary \tesserahash{}es encode to two canonical text forms. Both are fixed-length and round-trip byte-identical with the binary.

\subsection{Base64URL (RFC 4648 §5)}

73 bytes $\to$ 98 characters. The Base64URL alphabet \citep{rfc4648}: \code{A--Z}, \code{a--z}, \code{0--9}, \code{-}, \code{\_}; no padding (no trailing \code{=}). Encoding consumes 24 bits per 4 characters; 73 bytes = 584 bits = 24 full 3-byte groups (96 characters) plus 1 trailing byte (8 bits), and that trailing byte requires 2 more characters --- 98 in total. The final character carries 2 bits of value plus 4 bits of zero padding that MUST be zero (so the last character is always one of \code{A}, \code{Q}, \code{g}, \code{w}). The decoder must validate exactly 98 input characters.

\textbf{Decode validation.} A canonical decoder:
\begin{enumerate}
  \item Rejects input whose length is not 98 characters.
  \item Rejects any character outside the URL-safe alphabet (no \code{+}, \code{/}, \code{=}).
  \item Rejects inputs whose final character decodes to a 6-bit value with non-zero low 4 bits (equivalently: the last character must be one of \code{A}, \code{Q}, \code{g}, \code{w} --- the four 6-bit values whose low 4 bits are zero).
\end{enumerate}

The third check is the padding-bit guard. Without it, two distinct 98-character strings would decode to the same 73-byte binary, which violates canonical round-trip. The KAT file \texttt{tessera\_hash\_decode\_b64.json} pins this rejection.

\subsection{Crockford Base32 (canonical form)}

73 bytes $\to$ 117 characters. The Crockford Base32 alphabet \citep{crockford}: \code{0--9} then \code{A--Z} minus \code{I/L/O/U}. Encoding consumes 5 bits per character; 73 bytes = 584 bits = 116 full quintets plus 4 bits, so 117 characters carry the 4-bit excess in the final character with 1 bit of zero padding.

\textbf{Decode validation.}
\begin{enumerate}
  \item Length exactly 117 characters.
  \item Case-insensitive: lower-case folds to upper-case before lookup.
  \item Confusable canonicalization: \code{I/L $\to$ 1}, \code{O $\to$ 0} accepted on decode. The encoder emits canonical glyphs only.
  \item Padding bit: the final character's 5-bit value MUST have its low bit zero. Equivalent to the Base64URL padding-bit guard but in the Crockford alphabet.
\end{enumerate}

\textbf{Sort property.} Crockford Base32 ASCII sort preserves the byte order of the underlying 73-byte binary, which preserves generation-time order (because the leading 6 bytes are the big-endian timestamp). Base64URL ASCII sort \emph{does not} preserve byte order --- its alphabet ordering does not match ASCII codepoint order. Consumers that need a textual column to sort by generation time should use the Crockford form or compare the underlying bytes.

\subsection{Why two encodings}

\begin{itemize}
  \item \textbf{Base64URL} is the right choice for JSON values, HTTP headers, query strings, and database columns whose ordering does not matter. It is denser (98 chars vs 117) and is the canonical web encoding.
  \item \textbf{Crockford Base32} is the right choice for log lines, CLI tools, human transcription (the alphabet excludes the I/L/O/U confusables on purpose), and any database column where lexicographic sort must reflect generation time.
\end{itemize}

Both choices mirror TesseraID's encoding rationale exactly; the two projects use the same encoder code paths in every language (with different alphabet sizes per project).

\subsection{Examples}

\begin{itemize}
\item Binary \tesserahash{} of the empty string \texttt{b""} at timestamp $1700000000000$ and counter $0$:\\
73-byte binary $\to$ 98-char Base64URL $\to$ 117-char Crockford. Test vectors with byte-exact expected outputs ship in \texttt{tessera-hash/conformance/vectors/tessera\_hash\_encode\_b64.json} and \texttt{tessera\_hash\_encode\_crockford.json}.
\item Decoder round-trip: any 73-byte binary $b$ satisfies $\code{decode\_b64}(\code{encode\_b64}(b)) = b$ and $\code{decode\_b32}(\code{encode\_b32}(b)) = b$.
\end{itemize}
